% !TeX root = ../main.tex
% Original qualitative illustration: neither curve is a measured ground truth
% or a result of a proposed method. Both pass the same three sampled positions.
\begingroup
% Shared local styles for the four original problem illustrations in Section 2.2.
% Loaded inside each figure's group; all meanings remain legible in grayscale.
\linespread{1}\selectfont
\tikzset{
  rpBase/.style={x=1cm,y=1cm,text=black,line cap=round,line join=round,
    font=\songti\fontsize{9.8pt}{12pt}\selectfont,
    every node/.style={align=center,inner sep=0pt,outer sep=0pt}},
  rpTitle/.style={font=\heiti\fontsize{10.3pt}{13pt}\selectfont},
  rpNote/.style={font=\songti\fontsize{9.3pt}{12pt}\selectfont},
  rpFrame/.style={draw=black,line width=0.65pt},
  rpArrow/.style={draw=black,line width=0.8pt,
    -{Latex[length=1.7mm,width=1.15mm]}},
  rpMotion/.style={draw=black,line width=0.95pt,
    -{Latex[length=1.6mm,width=1.1mm]}},
  rpAction/.style={draw=black,line width=1.5pt,
    -{Latex[length=2.2mm,width=1.7mm]}},
  rpHistory/.style={draw=black!65,line width=0.65pt,
    dash pattern=on 1.4pt off 1.6pt},
  rpFuture/.style={draw=black,line width=0.75pt,
    dash pattern=on 3pt off 2pt},
  rpBoundary/.style={draw=black,line width=1.25pt}
}
\newcommand{\rpPanel}[5]{%
  \draw[rpFrame] (#1,#2) rectangle (#3,#4);
  \path[fill=black!5] (#1+0.01,#4-0.55) rectangle (#3-0.01,#4-0.01);
  \draw[rpFrame] (#1,#4-0.55) -- (#3,#4-0.55);
  \node[rpTitle] at ({(#1+#3)/2},{#4-0.275}) {#5};
}
\newcommand{\rpFooter}[1]{%
  \draw[rpFrame] (0.10,0.10) rectangle (15.90,0.77);
  \node at (8,0.435) {#1};
}
% Hatching is explicitly drawn and clipped, so no additional TikZ libraries
% or color meanings are required. Argument: a closed local-coordinate path.
\newcommand{\rpHatchFace}[1]{%
  \path[fill=black!5] #1;
  \begin{scope}
    \clip #1;
    \foreach \rhx in {-6,-5.8,...,8}{%
      \draw[black,line width=0.35pt] (\rhx,-1) -- ++(3,6);
    }
  \end{scope}
  \draw[rpBoundary] #1;
}

\begin{tikzpicture}[rpBase]
  \path[use as bounding box] (0,0) rectangle (16,5.35);
  \rpPanel{0.10}{1.05}{5.65}{5.20}{有限离散观测}
  \rpPanel{6.35}{1.05}{15.90}{5.20}{多种连续运动解释}

  % Three observations of the same object in the same image coordinate frame.
  \foreach \xl/\xc/\idx in {0.42/1.13/1,2.18/2.89/2,3.94/4.65/3}{%
    \draw[draw=black,line width=0.5pt] (\xl,2.50) rectangle ++(1.42,1.35);
    \draw[draw=black!45,line width=0.4pt] (\xl+0.08,2.82) -- ++(1.26,0);
    \node at (\xc,2.18) {$t_{\idx}$};
  }
  \filldraw[draw=black,fill=black!12,line width=0.7pt] (0.80,3.18) circle[radius=0.16];
  \filldraw[draw=black,fill=black!12,line width=0.7pt] (2.88,3.18) circle[radius=0.16];
  \filldraw[draw=black,fill=black!12,line width=0.7pt] (4.95,3.18) circle[radius=0.16];
  \draw[rpFrame] (0.42,1.30) rectangle (5.33,1.85);
  \node[rpNote] at (2.875,1.575) {采样位置已知，帧间过程待恢复};
  \draw[rpArrow] (5.73,3.05) -- (6.27,3.05);

  % Position against time. The continuous candidates agree at ALL observations.
  \draw[rpArrow,line width=0.55pt] (7.05,2.00) -- (15.35,2.00);
  \draw[rpArrow,line width=0.55pt] (7.05,2.00) -- (7.05,4.30);
  \node at (15.48,2.00) {$t$};
  \node[anchor=east] at (7.00,4.28) {$x(t)$};
  \foreach \xx/\yy/\idx in {7.55/2.50/1,10.80/3.10/2,14.05/3.70/3}{%
    \draw[draw=black!45,line width=0.4pt,densely dotted] (\xx,2.00) -- (\xx,\yy);
    \node[rpNote] at (\xx,1.77) {$t_{\idx}$};
  }
  \draw[draw=black,line width=0.95pt] (7.55,2.50) -- (14.05,3.70);
  \draw[draw=black,line width=0.9pt,dash pattern=on 3pt off 2pt]
    (7.55,2.50) .. controls (8.35,4.40) and (9.80,4.40) .. (10.80,3.10)
    .. controls (11.65,1.95) and (13.10,1.90) .. (14.05,3.70);
  \foreach \xx/\yy in {7.55/2.50,10.80/3.10,14.05/3.70}{%
    \filldraw[fill=black,draw=white,line width=0.5pt] (\xx,\yy) circle[radius=0.065];
  }
  \draw[draw=black,line width=0.95pt] (8.05,1.36) -- (8.85,1.36);
  \node[rpNote,anchor=west] at (9.05,1.36) {解释 A};
  \draw[draw=black,line width=0.9pt,dash pattern=on 3pt off 2pt] (11.55,1.36) -- (12.35,1.36);
  \node[rpNote,anchor=west] at (12.55,1.36) {解释 B};

  \rpFooter{研究关注：有效观测如何约束帧间运动？}
\end{tikzpicture}
\endgroup
